\input{../tex-inputs/latex-leseansicht-vorspann}

\section[Arthur und Olga Schnitzler an Gustav Schwarzkopf, 6. 11. 190{[}3{]}]{L04379 Arthur und Olga Schnitzler an Gustav Schwarzkopf, 6. 11. 190{[}3{]}}
\nopagebreak\mylabel{L04379v}
\rehead{ }\normalsize\beginnumbering\briefempfaengerindex{Schwarzkopf, Gustav@\textsc{Schwarzkopf, Gustav}!zzzSchnitzler, Olga@\emph{von Olga Schnitzler}!1903-11-061@{6. 11. 190{[}3{]}}|(be}\briefempfaengerindex{Schwarzkopf, Gustav@\textsc{Schwarzkopf, Gustav}!zzzSchnitzler, Arthur@\emph{von Arthur Schnitzler}!1903-11-061@{6. 11. 190{[}3{]}}|(be}
\toendnotes[C]{\smallbreak\pagebreak[2]}
\correspDesc{Versand  durch Arthur Schnitzler, Olga Schnitzler am 6. 11. 190[3] in Semmering
\newline{}Übermittlung  am 6. 11. 1903 in Semmering
\newline{}Zustellung  am 7. 11. 190[3] in Wien
\newline{}Erhalt  durch Gustav Schwarzkopf am 7. 11. 190[3] in Wien}\toendnotes[C]{\smallbreak}
\Standort{DLA, A:Schnitzler, HS.1985.1.1897.}
\physDesc{Bildpostkarte, 108 Zeichen
\newline{}Handschrift Arthur Schnitzler: Bleistift, deutsche Kurrent
\newline{}Handschrift Olga Schnitzler: Bleistift, lateinische Kurrent
\newline{}Versand: 1) Stempel: »\nobreak{}\oindex{Semmering@\textbf{Semmering}|pwk}Semmering 1, 6. 11. 03, 8\nobreak{}«.   2) Stempel: »\nobreak{}\oindex{I., Innere Stadt@\textbf{I., Innere Stadt}|pwk}Wien {[}1/1{]} 1, 7. 11. \textcolor{gray}{03}, 8–9½V., Bestellt\nobreak{}«. }\toendnotes[C]{\smallbreak}\pstart{}{\pb}Herrn Gustav Schwarzkopf\pend{}\pstart{}Wien I\oindex{I., Innere Stadt@\textbf{I., Innere Stadt}|pw}\pend{}\pstart{}Tiefer Graben 23\oindex{Wien@\textbf{Wien}!I., Innere Stadt@\textbf{I., Innere Stadt}!Tiefer Graben 23@\textbf{Tiefer Graben 23}|pw}.\pend{}{\bigskip}
\pstart
           \centering{}{\pb}\textcolor{gray}{\textbf{SEMMERING\oindex{Semmering@\textbf{Semmering}|pw} Partie am am Hochweg\oindex{Hochstraße [Semmering]@\textbf{Hochstraße [Semmering]}|pw}.}}\pend
           \vspace{1em}
\pstart
           \raggedleft{}{\pb}\label{K_L04371-1v}\edtext{6. 11. 90\textcolor{gray}{3}}{\lemma{\textnormal{\emph{6. 11. 903}}}\Cendnote{\textnormal{Obzwar die Jahresziffer am Poststempel auch
                     auf das Jahr 1908 verweisen könnte,
                     ist dies durch das Vorhandensein eines Empfangsstempels auszuschließen, da solche nur bis 1907 im Einsatz waren.}}}\label{K_L04371-1}.\pend
           \vspace{0.5em}
\pstart
           Herzliche Grüße!\pend
           \pstart Ihr \spacefill\mbox{Arthur Schn}\pend{}\selectlanguage{ngerman}\vspace{1em}
\pstart
           \noindent{}{[}hs. O. Schnitzler:{]} Auch von mir!\pend
           \pstart \spacefill\mbox{O. Sch.}\pend{}\selectlanguage{ngerman}\endnumbering\briefempfaengerindex{Schwarzkopf, Gustav@\textsc{Schwarzkopf, Gustav}!zzzSchnitzler, Olga@\emph{von Olga Schnitzler}!1903-11-061@{6. 11. 190{[}3{]}}|)be}\briefempfaengerindex{Schwarzkopf, Gustav@\textsc{Schwarzkopf, Gustav}!zzzSchnitzler, Arthur@\emph{von Arthur Schnitzler}!1903-11-061@{6. 11. 190{[}3{]}}|)be}\mylabel{L04379h}
\begin{anhang}
\end{anhang}\newcommand{\dateiname}{L04379}\newcommand{\titel}{Arthur und Olga Schnitzler an Gustav Schwarzkopf, 6. 11. 190[3]}\newcommand{\editorInnen}{Herausgegeben von Jahnke, SelmaMüller, Martin Anton}\input{../tex-inputs/latex-leseansicht-abspann}
